% 6-Apr-2019
% Modified source from
% <https://www.dcl.hpi.uni-potsdam.de/media/theses/>
%
\documentclass[parskip=half, fleqn]{scrbook} % <= Druckversion: "scrbook", Bildschirmversion: "scrreprt"
% meine Einstellung
\usepackage{content/praeambel-print}

% ABOUT
\newcommand{\typ}{Projekt}
\newcommand{\autor}{Jan Unger}
\newcommand{\titel}{notizenDummyUbuntu-v02} % wird per Powershell Script ersetzt
\newcommand{\untertitel}{Mitschrift}
\newcommand{\ort}{Wuppertal}
\newcommand{\datum}{\today}
\newcommand{\website}{https://www.ju1.eu}

% DOCUMENT
%\KOMAoption{draft}{true} % <= z.B. zum "Debuggen" der Overfull-Boxes
\bibliography{content/literatur}% meine Literatur
\bibliography{content/literatur-laufen}% meine Literatur

\begin{document}
% Einband
%\pagenumbering{alph}
%\ifisbook% ju -  25-Apr-2018  - coverpage.tex 
% Modified source from
% <https://www.dcl.hpi.uni-potsdam.de/media/theses/>
%
\begin{titlepage}
	\setlength{\evensidemargin}{0.5\evensidemargin+0.5\oddsidemargin}
	\setlength{\oddsidemargin}{\evensidemargin}

	\centering

	\raisebox{-0.5\height}{\includegraphics[width=0.95\textwidth]{content/titelbild-black.pdf}}

	\vspace*{4\baselineskip}
	{\usekomafont{title}\titel\par}
	\vspace*{\baselineskip}
	{\usekomafont{author}Autor: \autor}\par

	\vfill

\end{titlepage}
\fi
%\ifisbook\cleardoubleemptypage\fi

% (Haupt-)Titelseite, Zusammenfassung, ggf. Danksagung & Inhaltsverzeichnis
\pagenumbering{roman}
% ju - 25-Apr-18 - titelpage.tex 
% Modified source from
% <https://www.dcl.hpi.uni-potsdam.de/media/theses/>
%
\begin{titlepage}
	\setlength{\evensidemargin}{0.8\evensidemargin+0.5\oddsidemargin}
	\setlength{\oddsidemargin}{\evensidemargin}

	\centering

	\raisebox{-0.5\height}{\includegraphics[width=8cm]{img/titelbild.pdf}}
	%\raisebox{-0.5\height}{\includegraphics[width=8cm]{img/titelbild.pdf}}
	%\hspace*{.2\textwidth}
	%\raisebox{-0.5\height}{\includegraphics[width=3cm]{img/logo.pdf}}
	\hfill
	\vspace{3mm}
	{\color{meinblue} \rule{\textwidth}{3pt}}


	\vspace*{4\baselineskip}
	%{\usekomafont{subject}\typ}\par
	%\vspace*{\baselineskip}
	{\usekomafont{title}\titel\par}
	\vspace*{\baselineskip}
	{\usekomafont{author}Autor: \autor}\par
	%\vspace*{\baselineskip}
	%{\usekomafont{date}\ort}\par


	\begin{flushleft}
		\vfill
		{\usekomafont{date}Version: \datum}\par
		\vspace*{\baselineskip}
		{\qrcode[hyperlink,height=1.5cm]{\website}}\par
	\end{flushleft}

	\setcounter{page}{1}

\end{titlepage}

%\ifisbook\cleardoubleemptypage\fi\include{content/zusammenfassung}% example
%\ifisbook\cleardoubleemptypage\fi\include{content/danksagung}% example
\tableofcontents
\cleardoublepage

% Textteil
\pagenumbering{arabic}

% Inhalt
% --------------------------------------------------
%
% ---------------------------------
\chapter{neu}
%----------
% \section{ }
% \subsection{ }\label{ }\index{ }
%----------

\hypertarget{test}{%
\section{Test}\label{test}}

\hypertarget{quellcode}{%
\subsection{Quellcode}\label{quellcode}}

(\autoref{code:Codename } Codename). % Codeverweis =

% Quellcode
\lstset{language=Bash} % C, [LaTeX]TeX, Bash, Python
\begin{lstlisting}[numbers=left, frame=l, framerule=0.1pt,%
% -----------
	caption={Codename }, % Caption
	label={code:Codename }  % Label
]% ---------=

  # Überschrift
  ## Überschrift 2
  ### Überschrift 3
\end{lstlisting}

\hypertarget{bild}{%
\subsection{Bild}\label{bild}}

% Bild Referenz
(\autoref{pic:logo} logo.pdf ). % Bildverweis = logo.pdf

\begin{figure}[!hp]% hier
	\centering
	\includegraphics[width=0.8\textwidth]{img/logo.pdf}
	% ---------
	\caption[logo.pdf ]{logo.pdf } % Caption
	\label{pic:logo}	% Referenz
	% ---------
\end{figure}

\hypertarget{tabelle}{%
\subsection{Tabelle}\label{tabelle}}

(\autoref{tab:Tabellenname } Tabellenname ). % Tabellenverweis = Tabellenname

% Tabelle
\begin{table}[!hp] % hier
	\centering
	%\setlength{\tabcolsep}{5mm} % Spaltenlänge fest
	\rowcolors{1}{}{lightgray!20}  % Farbe
	%\begin{tabularx}{\textwidth}{XX} % auto. Spaltenumbruch
	\begin{tabular} []{rll}
	\toprule % Inhalt
	% ---------
\textbf{Nr.} & \textbf{Begriffe} & \textbf{Erklärung}\\
	\midrule
1 & a1 & a2\\
2 & b1 & b2\\
3 & c1 & c2\\
	\bottomrule
	%\end{tabularx}
	\end{tabular}
	\caption[Tabellenname ]{Tabellenname }	% Caption
	\label{tab:Tabellenname }	% Referenz
	%-----------
\end{table}
%-----------



 % wird per Powershell Script ersetzt
%
% --------------------------------------------------
%\chapter{Kapitel}
%\chapter{Readme}
%\include{content/Readme}
%\chapter{MarkdownSpickzettel}
%%============================
% \section{ }
% \subsection{ }\label{ }\index{ }
%============================

\section{Markdown - Spickzettel}\label{markdown-syntax}

\subsection{Überschrift}\label{ueberschrift}

% Quellcode
\lstset{language=Bash} % C, [LaTeX]TeX, Bash, Python
\begin{lstlisting}[numbers=left, frame=l, framerule=0.1pt,%
% ======================
	caption={ }, % Caption
	label={code: }  % Label
]% =====================

    # Überschrift
    ## Überschrift 2
    ### Überschrift 3
\end{lstlisting}

\subsection{Bild}\label{bild}

% Bild Referenz
(\autoref{pic:logo1} logo.pdf). % Bildverweis = logo.pdf

\begin{figure}[!hp]% hier
	\centering
	\includegraphics[width=0.8\textwidth]{img/logo.pdf}
	% ====================
	\caption[logo.pdf ]{logo.pdf } % Caption
	\label{pic:logo1}	% Referenz
	% ====================
\end{figure}

% Quellcode
\lstset{language=Bash} % C, [LaTeX]TeX, Bash, Python
\begin{lstlisting}[numbers=left, frame=l, framerule=0.1pt,%
% ======================
	caption={ }, % Caption
	label={code: }  % Label
]% =====================

    # bild
    ![Logo](img/logo.pdf)
\end{lstlisting}

\subsection{Tabelle}\label{tabelle}

(\autoref{tab:Tabelle } ). % Tabellenverweis = Tabelle

% Tabelle
\begin{table}[!hp] % hier
	\centering
	%\setlength{\tabcolsep}{5mm} % Spaltenlänge fest
	\rowcolors{1}{}{lightgray!20}  % Farbe
	%\begin{tabularx}{\textwidth}{XX} % auto. Spaltenumbruch
	\begin{tabular} []{rll}
	\toprule % Inhalt
	% ---------
\textbf{Nr.} & \textbf{Begriffe} & \textbf{Erklärung}\\
	\midrule
1 & a1 & a2\\
2 & b1 & b2\\
3 & c1 & c2\\
4 & a1 & a2\\
	\bottomrule
	%\end{tabularx}
	\end{tabular}
	\caption[Tabelle ]{Tabelle }	% Caption
	\label{tab:Tabelle }	% Referenz
	%============================
\end{table}

% Quellcode
\lstset{language=Bash} % C, [LaTeX]TeX, Bash, Python
\begin{lstlisting}[numbers=left, frame=l, framerule=0.1pt,%
% ======================
	caption={ }, % Caption
	label={code: }  % Label
]% =====================

    # tabelle
    |**Nr.**|**Begriffe**|**Erklärung**|
    |--:|:-----|:----|
    | 1     | a1         | a2                |
    | 2     | b1         | b2              |
    | 3     | c1         | c2              |
    | 4     | a1         | a2              |
\end{lstlisting}

\subsection{Mathe}\label{mathe}

$[ V ] = [ \Omega ] \cdot [ A ]$ o. $U = R \cdot I$ o.
$R = \frac{U}{I}$

% Quellcode
\lstset{language=Bash} % C, [LaTeX]TeX, Bash, Python
\begin{lstlisting}[numbers=left, frame=l, framerule=0.1pt,%
% ======================
	caption={ }, % Caption
	label={code: }  % Label
]% =====================

    # Mathe
    $[ V ] = [ \Omega ] \cdot [ A ]$ o. $U = R \cdot I$ o. $R = \frac{U}{I}$
\end{lstlisting}

\textbf{Matheumgebung:}

\begin{align*}
    \sum_{i=1}^5 a_i = a_1 + a_2 + a_3 + a_4 + a_5
\end{align*}

% Quellcode
\lstset{language=Bash} % C, [LaTeX]TeX, Bash, Python
\begin{lstlisting}[numbers=left, frame=l, framerule=0.1pt,%
% ======================
	caption={ }, % Caption
	label={code: }  % Label
]% =====================

    # Matheumgebung
    \begin{align*}
        \sum_{i=1}^5 a_i = a_1 + a_2 + a_3 + a_4 + a_5
    \end{align*}
\end{lstlisting}

\subsection{Absätze}\label{absaetze}

Dies hier ist ein Blindtext zum Testen von Textausgaben. Wer diesen Text
liest, ist selbst schuld. Der Text gibt lediglich den Grauwert der
Schrift an. Ist das wirklich so? Ist es gleichgültig, ob ich schreibe:
\frqq Dies ist ein Blindtext\flqq\, oder \frqq Huardest gefburn\flqq\,? Kjift -
mitnichten! Ein Blindtext bietet mir wichtige Informationen.

Dies hier ist ein Blindtext zum Testen von Textausgaben. Wer diesen Text
liest, ist selbst schuld. Der Text gibt lediglich den Grauwert der
Schrift an. Ist das wirklich so? Ist es gleichgültig, ob ich schreibe:
\frqq Dies ist ein Blindtext\flqq\, oder \frqq Huardest gefburn\flqq\,? Kjift -
mitnichten! Ein Blindtext bietet mir wichtige Informationen.

\subsection{Listen}\label{listen}

\textbf{ungeordnete Liste}

\begin{itemize}% Liste Punkt
\item
  a
\item
  b

  \begin{itemize}% Liste Punkt
  \item
    bb
  \end{itemize}
\item
  c
\end{itemize}

% Quellcode
\lstset{language=Bash} % C, [LaTeX]TeX, Bash, Python
\begin{lstlisting}[numbers=left, frame=l, framerule=0.1pt,%
% ======================
	caption={ }, % Caption
	label={code: }  % Label
]% =====================

    # ungeordnete Liste
    - a
    - b
        - bb
    - c
\end{lstlisting}

\textbf{Sortierte Liste}

\begin{enumerate}% Liste 1) oder a)
\item
  eins
\item
  zwei
\item
  drei
\end{enumerate}

% Quellcode
\lstset{language=Bash} % C, [LaTeX]TeX, Bash, Python
\begin{lstlisting}[numbers=left, frame=l, framerule=0.1pt,%
% ======================
	caption={ }, % Caption
	label={code: }  % Label
]% =====================

    # Sortierte Liste
    1. eins
    2. zwei
    3. drei
\end{lstlisting}

\textbf{Sortierte Liste}

\begin{enumerate}% Liste 1) oder a)
\def\labelenumi{\alph{enumi})}
\item
  a
\item
  b
\item
  c
\end{enumerate}

(\autoref{code:Liste } Liste). % Codeverweis = Liste

% Quellcode
\lstset{language=Bash} % C, [LaTeX]TeX, Bash, Python
\begin{lstlisting}[numbers=left, frame=l, framerule=0.1pt,%
% ======================
	caption={Liste }, % Caption
	label={code:Liste }  % Label
]% =====================

    # Sortierte Liste
    a) a
    b) b
    c) c
\end{lstlisting}

\subsection{Anführungszeichen}\label{anfuehrungszeichen}

\frqq Anführungszeichen\flqq\, oder \flqq Anführungszeichen\frqq   oder
\frqq Anführungszeichen\flqq

% Quellcode
\lstset{language=Bash} % C, [LaTeX]TeX, Bash, Python
\begin{lstlisting}[numbers=left, frame=l, framerule=0.1pt,%
% ======================
	caption={ }, % Caption
	label={code: }  % Label
]% =====================

    # Anführungszeichen
    "Anführungszeichen" oder
    \flqq Anführungszeichen\frqq   oder
    \frqq Anführungszeichen\flqq
\end{lstlisting}

\subsection{Quellenangabe}\label{quellenangabe}

Quelle \textcite{schlosser_latex:2016}

Quelle Text\footfullcite{schlosser_latex:2016}

Quelle \cite{schlosser_latex:2016}

% Quellcode
\lstset{language=Bash} % C, [LaTeX]TeX, Bash, Python
\begin{lstlisting}[numbers=left, frame=l, framerule=0.1pt,%
% ======================
	caption={ }, % Caption
	label={code: }  % Label
]% =====================

    # Quellenangabe
    Quelle \textcite{schlosser_latex:2016}
    Quelle Text\footfullcite{schlosser_latex:2016}
    Quelle \cite{schlosser_latex:2016}
\end{lstlisting}
%============================
%============================
%============================

%\chapter{LatexSpickzettel}
%\include{content/LatexSpickzettel}
%\chapter{AdobeStockLizenz}
%%----------
% \section{ }
% \subsection{ }\label{ }\index{ }
%----------

%----
% \section{ }
% \subsection{ }\label{ }\index{ }
%----

\section{Einschränkungen der Standardlizenz von Adobe
Stock}\label{einschraenkungen-der-standardlizenz-von-adobe-stock}

Unter einer Standardlizenz zulässig:

\begin{itemize}% Liste Punkt% Liste Punkt
\item
  Reproduktion von bis zu 500.000 Kopien des Stockmediums auf
  Produktverpackungen und in gedruckten Marketingmaterialien sowie in
  digitalen Dokumenten oder Software.
\item
  Einschließen des Stockmediums in E-Mail-Marketing, mobile Werbung oder
  Fernsehprogramme, wenn weniger als 500.000 Zuschauer zu erwarten sind.
\item
  Veröffentlichen des Stockmediums auf einer Website ohne Einschränkung
  der Anzahl an Besucher. Wenn das Stockmedium unverändert auf einer
  Website der sozialen Medien veröffentlicht wird, ist eine
  Namensnennung erforderlich ((c) Autorenname -- stock.adobe.com).
\item
  Einschließen des Stockmediums in Produkte auf geringfügige Weise, z.
  B. in Lehrbüchern.
\end{itemize}

Unter einer Standardlizenz nicht zulässig:

Erstellen von Waren oder Handels- und Vertriebsprodukten, bei denen das
Stockmedium selbst den Hauptwert des Produkts ausmacht. So dürfen Sie
beispielsweise keine Poster, T-Shirts oder Kaffeetassen gestalten, die
hauptsächlich wegen des aufgedruckten Stockmediums gekauft werden.

Bilder

((C) Autorenname -- stock.adobe.com).

((C) Sergey Nivens -- stock.adobe.com).
%-----
%-----------

%\chapter{VorlageLaTeX}
%%----------
% \section{ }
% \subsection{ }\label{ }\index{ }
%----------

%----
% \section{ }
% \subsection{ }\label{ }\index{ }
%----

\section{Vorlage (LaTeX)}\label{vorlage-latex}

Modified source from v.14

\url{https://www.dcl.hpi.uni-potsdam.de/media/theses/}

\subsection{Literaturverzeichnis, Zitate und
Quellenangaben}\label{literaturverzeichnis-zitate-und-quellenangaben}

Tool \url{https://www.zotero.org/}

\subsection{Rechtschreibung \&
Grammatik}\label{rechtschreibung-grammatik}

\url{https://www.duden.de/Liste-der-rechtschreiblich-schwierigen-Woerter}

\subsection{Drucken \& Binden}\label{drucken-binden}

\url{http://www.potsdam-druck.de/digitaldruck/softcoverbuecher-drucken-binden.html}
%-----
%-----------

%
%\input{Projekt-files}
% --------------------------------------------------
%
% ggf. Anhang
%\appendix\include{content/anhang} % example

% Bibliographie
\ifisbook\cleardoubleemptypage\fi
\phantomsection\addcontentsline{toc}{chapter}{\refname}
\printbibliography[category=cited]
%\printbibliography[heading=bibintoc] % Bibliographie

% Eigenstaendigkeitserklaerung
%\ifisbook\pagestyle{plain}\cleardoubleemptypage\include{content/erklaerung}\fi % example

\end{document}
%============================
%============================
%============================
